% ===================================================================
%  TABEL Nr. 10 — Die see. — Die hawe.
%  Meme regime que les precedents. Tableau a UNE SEULE COLONNE : les
%  cles vont t10-01-1 a t10-25-2, sans c1 ni c2, et l'apparat porte
%  le titre et le premier intertitre ensemble — il n'y a pas de
%  t10-tit-1.
%
%  UNE PHRASE ENTIERE EST EN GRAS AU 5, dans le fac-simile ido comme
%  dans le neerlandais : ce n'est pas un terme de planche, c'est le
%  compositeur qui a laisse courir son gras. On le reproduit,
%  ponctuation comprise, parce que les autres colonnes le
%  reproduisent.
%
%  LE TABLEAU DISTINGUE DEUX SORTES DE DOCK : « flotacanta doki »
%  (27) sont les docks flottants ou l'on repare, « flot-doko » (53)
%  le bassin a flot ou les barques s'abritent. L'afrikaans les
%  separe aussi — drywende dokke et vlotdok.
%
%  ET C'EST LE TABLEAU OU LE NEERLANDAIS PESE LE PLUS LOURD sur une
%  colonne afrikaans : presque tout le vocabulaire de la mer est
%  commun aux deux langues, si bien qu'on ecrit sans y penser
%  « zeilschepen », « scheepswerven », « schroef ». Ce sont
%  seilskepe, skeepswerwe, skroef. Le controle de forme lit le sch-
%  initial et le -ij- ; il ne lit pas ce qui se prononce pareil.
% ===================================================================

%%K t10-apar-1 apar
\VUsaut{4.0mm}
\VUtitre{15.0pt}{60}{TABEL Nr. 10}
\VUsaut{3.0mm}
\VUpk{11.4pt}{40}{Die see. --- Die hawe.}
\VUsaut{3.0mm}

%%K t10-01-1 p
1. --- In die somer, terwyl sommige rus en koelte op 'n berg soek, gaan ander liewer na
die kus. So het ons verlede jaar gedoen, want ons hou baie van die \VUgras{Oseaan}
\textsuperscript{(1)}.

%%K t10-02-1 p
2. --- Wat my betref, ek voel 'n groot genot wanneer ek daardie onmeetlike watervlakte
aanskou wat hom tot by die \VUgras{horison} \textsuperscript{(2)} uitstrek, en wanneer ek
die enorme \VUgras{stoomskepe} \textsuperscript{(3)} vir 'n lang \VUgras{oortog} sien
vertrek waarop die reisigers inskeep wat na Amerika gaan.

%%K t10-03-1 p
3. --- Daarom kies my vader, wat my neiging ken, om die vakansie deur te bring altyd 'n
baie rustige strand wat naby een van ons groot \VUgras{oorlogshawens} lê.

%%K t10-03-2 p
Tydens ons laaste verblyf het die \VUgras{eskader} agter die enorme \VUgras{seewering}
\textsuperscript{(4)} kom anker gooi wat die \VUgras{oorlogshawe} \textsuperscript{(5)}
afsluit en beskerm, en dit het my na wens geluk om die verskillende
\VUgras{oorlogskepe} te bewonder, die klein \VUgras{duikboot} \textsuperscript{(10)} wat
duik en onder die water verdwyn en ook die enorme \VUgras{pantserskip}
\textsuperscript{(9)} wat tot nou toe byna net die aanvalle van die \VUgras{torpedoboot}
\textsuperscript{(8)} gevrees het.

%%K t10-tit-2 sub
\VUsaut{3.0mm}
\VUpk{10.2pt}{40}{Die klein skepe.}
\VUsaut{3.0mm}

%%K t10-04-1 p
4. --- Laasgenoemde het al baie \VUgras{handig} die swak punt van die dik metaalpantser
weet te vind om daar die gevaarlike \VUgras{torpedo} in te steek, waarvan die ontploffing
die ondergang van die skip veroorsaak, maar tog was dit minder te vrees as die duikboot,
want die torpedojaers agtervolg dit maklik.

%%K t10-05-1 p
5. --- Ek kon ook die vinnige gepantserde \VUgras{kruisers} \textsuperscript{(6)} sien,
byna net so onkwesbaar soos die egte pantserskip, etlike \VUgras{transportskepe}
\textsuperscript{(12)} drie of vier \VUgras{kanonneerbote} \textsuperscript{(7)} en ten
slotte 'n \VUgras{fregat} \textsuperscript{(11)}. \VUgras{Die skouspel van daardie vloot,
wanneer dit manoeuvreer om die anker te lig, is die boeiendste.}

%%K t10-06-1 p
6. --- Watter verskil in bou tussen daardie groot staalmassas en die elegante
\VUgras{driemasters} of die sierlike \VUgras{seilskepe} \textsuperscript{(13)} wat nou
nog in die koopvaardyvloot gebruik word en wat ek met al hulle seile die hawe sien
binnevaar het toe ek op die \VUgras{kaai} \textsuperscript{(14)} gewandel het. Waarlik,
hulle het baie van hulle voordele verloor \VUgras{wanneer} die \VUgras{wind} teen was.
Hulle moes al die seile stryk om weer die bassin binne te vaar en 'n \VUgras{sleepboot}
\textsuperscript{(16)} was nodig om hulle te gaan haal en hulle, soos eenvoudige
\VUgras{vragskuite} \textsuperscript{(17)}, na die kaai te bring.

%%K t10-tit-3 sub
\VUsaut{3.0mm}
\VUpk{10.2pt}{40}{Die handelshawe.}
\VUsaut{3.0mm}

%%K t10-07-1 p
7. --- Wat boeiend is aan die hawe waaroor ek praat, is dat dit nie net 'n oorlogshawe
bevat nie, kunsmatig deur reusagtige werke aangelê, maar ook 'n wonderlike
\VUgras{handelshawe} wat in die \VUgras{monding} van 'n groot rivier lê.

%%K t10-08-1 p
8. --- Die landtong wat die twee bassins skei, is versterk en verdedig deur 'n reeks
\VUgras{batterye} en \VUgras{bastions} wat in die see vooruitspring. 'n Mens het 'n
besondere verlof nodig om op die \VUgras{walle} \textsuperscript{(15)} te wandel. Ek kon
dit maklik kry, deur my oom, wat ingenieur van die \VUgras{skeepswerwe}
\textsuperscript{(18)} is, en ek het die skepe gaan besoek wat 'n mens onder wye loodse
bou.

%%K t10-09-1 p
9. --- Ek het selfs die tewaterlating van 'n pantserskip bygewoon en ek onthou die
aangrypende oomblik wat die hele skare gevoel het toe die reusagtige massa, aan homself
oorgelaat, oor sy wiegvormige raamwerk begin gly het en op die water van die rivier kom
dryf het.

%%K t10-10-1 p
10. --- Om na die skeepswerwe te gaan, steek 'n mens die \VUgras{oefenterrein}
\textsuperscript{(19)} oor waar die soldate van die garnisoen elke dag kom oefen, 'n mens
gaan oor 'n yster \VUgras{bruggie} \textsuperscript{(20)} wat die paradeplein met die
\VUgras{arsenaal} \textsuperscript{(21)} verbind.

%%K t10-tit-4 sub
\VUsaut{3.0mm}
\VUpk{10.2pt}{40}{Die arsenaal en die dokke.}
\VUsaut{3.0mm}

%%K t10-11-1 p
11. --- Saam met my oom het ek daardie groot gebou besoek vol \VUgras{kanonne}
\textsuperscript{(22)}, \VUgras{koeëls} \textsuperscript{(23)} en \VUgras{granate}. Ek
het ook die \VUgras{metaalfabriek} \textsuperscript{(24)} gesien waarin die
\VUgras{ketels} \textsuperscript{(25)} van die stoomskepe vervaardig word.

%%K t10-12-1 p
12. --- Baie naby is die \VUgras{dokke} \textsuperscript{(26)} --- herstelddokke --- en
die baie merkwaardige \VUgras{drywende dokke} \textsuperscript{(27)} waarin ek van baie
naby, op 'n skip wat herstel word, die \VUgras{skroef} \textsuperscript{(28)}, die
\VUgras{kiel} \textsuperscript{(29)} en die \VUgras{roer} \textsuperscript{(30)} van 'n
boot kon sien.

%%K t10-13-1 p
13. --- Die handelshawe is net so bestudeerwaardig soos die oorlogshawe, en ek het baie
ure deurgebring met gaap op die kaaie waar die \VUgras{radbote} \textsuperscript{(31)}
gemeer lê wat die rivier weer opvaar, en met kyk hoe die \VUgras{mashyskrane}
\textsuperscript{(32)} werk, wat soos vere enorme vragte oplig.

%%K t10-tit-5 sub
\VUsaut{4.0mm}
\VUpk{10.2pt}{40}{Die loods.}
\VUsaut{3.0mm}

%%K t10-14-1 p
14. --- Ek het die \VUgras{transatlantiese skepe} \textsuperscript{(34)} bewonder wat die
ree uitvaar, met hulle \VUgras{vlae} \textsuperscript{(35)} in die volle wind, gelei deur
'n handige \VUgras{loods} \textsuperscript{(36)} wat wonderlik weet om die
\VUgras{vaargeultjie} te volg wat met \VUgras{boeie} \textsuperscript{(40)} en
\VUgras{bakens} gemerk is, terwyl hy die \VUgras{ligters} \textsuperscript{(41)} vermy
wat hulle moeisaam met hulle enigste vierkantige seil laat stuur. Ek het op die
\VUgras{voorskip} \textsuperscript{(37)} --- die boeg --- die \VUgras{hutte}
\textsuperscript{(39)} van die tweede klas onderskei, en op die \VUgras{agterskip}
\textsuperscript{(38)} --- die spieël --- die salonne vir die passasiers van die eerste
klas.

%%K t10-15-1 p
15. --- Soms het ek ook die laai bygewoon van 'n \VUgras{vragskuit}
\textsuperscript{(42)} waarin die ware van die \VUgras{pakhuis} \textsuperscript{(43)} en
van die \VUgras{seestasie} \textsuperscript{(44)} pas opgestapel is, aangevoer deur die
talle treine van die \VUgras{kaailyn} \textsuperscript{(45)}.

%%K t10-tit-6 sub
\VUsaut{3.0mm}
\VUpk{10.2pt}{40}{Roei vir plesier.}
\VUsaut{3.0mm}

%%K t10-16-1 p
16. --- Dikwels het ek in die \VUgras{vlotdok} \textsuperscript{(53)} geroei, waarin die
\VUgras{bote} \textsuperscript{(46)} kom skuil, en dit was 'n vreugde vir my om die
\VUgras{roeispaan} \textsuperscript{(47)} te hanteer. Ek het op die \VUgras{dek}
\textsuperscript{(48)} van 'n \VUgras{seilboot} \textsuperscript{(49)} geklim, ek het
langs \VUgras{toue} \textsuperscript{(52)} in die \VUgras{mas} \textsuperscript{(51)} tot
in die \VUgras{ra's} \textsuperscript{(50)} geklouter, daarna het ek my tog
\VUgras{met 'n jol} \textsuperscript{(54)} herhaal tussen swaar \VUgras{vissersbote}
\textsuperscript{(55)}, waar die \VUgras{nette} \textsuperscript{(56)} droog word.

%%K t10-17-1 p
17. --- Op die \VUgras{kaai} \textsuperscript{(58)} het die mees uiteenlopende
\VUgras{ware} \textsuperscript{(59)} voor my verbygetrek, in \VUgras{kiste}
\textsuperscript{(60)} of in \VUgras{bale} \textsuperscript{(61)}. Hulle het die
\VUgras{vrag} \textsuperscript{(63)} van die \VUgras{ligters} gevorm, en kragtige
\VUgras{krane} \textsuperscript{(62)} het hulle afgehaal en op \VUgras{vragwaens}
\textsuperscript{(64)} neergesit terwyl die \VUgras{vragryers} hulle aangegee en die
regte in die \VUgras{doeanekantoor} \textsuperscript{(65)} betaal het.

%%K t10-18-1 p
18. --- Ten slotte, toe ek by die \VUgras{draaibrug} \textsuperscript{(66)} of by die
\VUgras{sluis} \textsuperscript{(69)} gekom het, terwyl ek verby die
\VUgras{baggermasjien} \textsuperscript{(68)} gegaan het wat die \VUgras{kanaal}
\textsuperscript{(67)} skoonmaak, het ek op die kaai naby die \VUgras{seinpaal}
\textsuperscript{(70)} aan wal gegaan.

%%K t10-19-1 p
19. --- Aan die ander kant van daardie kaai kom juis die \VUgras{sloepe}
\textsuperscript{(71)} aanlê wat die offisiere van die eskader aan land bring. Dit is 'n
bewonderenswaardige skouspel om die matrose in die maat hulle roeispane te sien oplig,
terwyl die boot wat deur die \VUgras{golf} \textsuperscript{(72)} gedra word, maklik by
die trap van die aanlêplek kom. Op daardie plek kan 'n mens die \VUgras{gety} en die
beweging van \VUgras{vloed} en \VUgras{eb} op die \VUgras{strand} \textsuperscript{(73)}
beter waarneem.

%%K t10-tit-7 sub
\VUsaut{3.0mm}
\VUpk{10.2pt}{40}{Die offisiersklub.}
\VUsaut{3.0mm}

%%K t10-20-1 p
20. --- Om huis toe terug te keer het ek die \VUgras{elektriese trem}
\textsuperscript{(74)} gebruik waarvan die \VUgras{halte} \textsuperscript{(75)} naby die
\VUgras{paal} is wat voor die klub van die offisiere staan.

%%K t10-20-2 p
Van die \VUgras{balkon} \textsuperscript{(76)} van die klub af, versier met
\VUgras{ankers} \textsuperscript{(77)}, kanonne en koeëls, het ek dikwels 'n skitterende
versameling seeoffisiere gesien: \VUgras{luitenante} \textsuperscript{(78)} met hulle
degen, \VUgras{kapteins van die artillerie} \textsuperscript{(79)} en van die
\VUgras{infanterie} \textsuperscript{(80)} met hulle \VUgras{sabel}. Agter hulle het 'n
\VUgras{matroos} \textsuperscript{(81)} gestaan wat vir hulle 'n \VUgras{verkyker}
gebring het wanneer hulle wou onderskei wat in die verte gebeur.

%%K t10-21-1 p
21. --- Ek het ook jong \VUgras{onderluitenante} \textsuperscript{(82)} gesien wat die
bevele van hulle \VUgras{kommandant} \textsuperscript{(83)} of van die \VUgras{admiraal}
\textsuperscript{(84)} ontvang het, terwyl 'n \VUgras{kwartiermeester}
\textsuperscript{(85)} gewag het om hulle na die skip te bring.

%%K t10-22-1 p
22. --- Van daardie balkon af is die uitsig pragtig: 'n mens oorsien die hele strand met
sy \VUgras{tente} \textsuperscript{(86)} en \VUgras{badhuisies} \textsuperscript{(87)},
en 'n mens sien, op die uur van die bad, die \VUgras{baaiers} \textsuperscript{(88)} wat
vrolik dartel voor die \VUgras{casino} \textsuperscript{(89)}.

%%K t10-23-1 p
23. --- Aan die uiteinde van die strand vorm die kus 'n klein rotsagtige
\VUgras{skiereiland} \textsuperscript{(91)}, waar die \VUgras{brandergolf}
\textsuperscript{(92)} hom kom breek en waarop 'n \VUgras{vuurtoring}
\textsuperscript{(93)} gebou is wat snags die weg aan die skeepvaarders wys. Daardie
skiereiland is met die land net deur 'n baie nou \VUgras{landengte}
\textsuperscript{(90)} verbind.

%%K t10-tit-8 sub
\VUsaut{3.0mm}
\VUpk{10.2pt}{40}{Algemene gesig op die kuste.}
\VUsaut{3.0mm}

%%K t10-24-1 p
24. --- Die \VUgras{steilkus} \textsuperscript{(94)} strek hom uit, oopgeskeur deur die
see, wat daarin grillig 'n reeks \VUgras{baaie} \textsuperscript{(95)} en
\VUgras{voorgebergtes} (kape) teken, wat hom hoogs skilderagtig maak.

%%K t10-24-2 p
Op die \VUgras{plato} \textsuperscript{(96)}, wat die \VUgras{(see)engte}
\textsuperscript{(98)} oorheers, staan 'n baie belangrike \VUgras{fort}
\textsuperscript{(97)} en enkele « villa's ».

%%K t10-25-1 p
25. --- Daardie engte skei van die vasteland 'n baie gevaarlike \VUgras{eiland}
\textsuperscript{(99)}, omring deur \VUgras{riwwe} \textsuperscript{(100)} en
\VUgras{branders}, waar bote en selfs groot skepe soms strand. Ondanks die gereedheid van
die hulp en die vinnige aankoms van die \VUgras{reddingsbote}, is daardie
\VUgras{skipbreuke} verskriklik, en tydens die \VUgras{storms} van die nagewening het
etlike skepe met mense en ware vergaan.

%%K t10-25-2 p
Ondanks daardie nabyheid word \VUgras{die hawe} baie deur die seelui besoek.
\VUsaut{6.0mm}
\VUfilet{22mm}
